\documentclass[
    12pt,           % tamanho da fonte do texto
    oneside,        % impressão em um lado da folha
    a4paper,        % tamanho do papel
    chapter=TITLE,  % títulos de capítulos em caixa alta
    english,        % idioma adicional para pacotes internos
    brazil          % idioma principal do documento
]{abntex2}

\usepackage[utf8]{inputenc}      % codificação de entrada UTF-8
\usepackage[T1]{fontenc}         % codificação de fonte
\usepackage{indentfirst}         % recuo na primeira linha dos parágrafos
\usepackage{graphicx}            % inserção de imagens (uso futuro)
\usepackage{microtype}           % melhorias tipográficas sutis
\usepackage{newtxtext}           % fonte Times New Roman
\usepackage{pifont}              % ícones para tabela (check e cross)
\usepackage{xcolor}              % suporte a cores

% MARGENS superior 3 cm | esquerda 3 cm | inferior 2 cm | direita 2 cm
\usepackage[
    top=3cm,
    left=3cm,
    bottom=2cm,
    right=2cm
]{geometry}


% ESPAÇAMENTO E FORMATAÇÃO DE PARÁGRAFO
\OnehalfSpacing                  % espaçamento 1,5 entre linhas
\setlength{\parindent}{1.25cm}   % recuo de parágrafo
\setlength{\parskip}{0pt}        % sem espaçamento extra entre parágrafos


% FONTES E TAMANHOS DOS TÍTULOS (CAPÍTULOS, SEÇÕES E SUBSEÇÕES)
\renewcommand{\ABNTEXchapterfontsize}{\fontsize{12}{14}\selectfont}
\renewcommand{\ABNTEXsectionfontsize}{\fontsize{12}{14}\selectfont}
\renewcommand{\ABNTEXsubsectionfontsize}{\fontsize{12}{14}\selectfont}
\renewcommand{\ABNTEXsubsubsectionfontsize}{\fontsize{12}{14}\selectfont}

\renewcommand{\ABNTEXchapterfont}{\rmfamily\bfseries}
\renewcommand{\ABNTEXsectionfont}{\rmfamily\bfseries}
\renewcommand{\ABNTEXsubsectionfont}{\rmfamily\bfseries}
\renewcommand{\ABNTEXsubsubsectionfont}{\rmfamily\itshape}

\setlength{\beforechapskip}{0pt}
\setlength{\afterchapskip}{\baselineskip}


% TABELAS
\newcommand{\cmark}{\textcolor{green!60!black}{\ding{51}}}  % Símbolo de visto (verde)
\newcommand{\xmark}{\textcolor{red}{\textbf{X}}}            % Símbolo de X (vermelho)


% INFORMAÇÕES DO DOCUMENTO
\titulo{Não temos titulo ainda pois não sabemos oq fazer e blablabla blablabla blablabla blablabla blablabla blablabla blablabla}
\autor{%
    Eduarda Luiza Pinheiro Nepomuceno \\
    Jean Yuki Kimura \\
    Jhenifer Laís Barbosa \\
    João Pedro Duarte Giatti
}

% CONFIGURAÇÕES DE PDF (metadados)
\makeatletter
\hypersetup{
    pdftitle={\@title},
    pdfauthor={\@author},
    pdfsubject={Engenharia de Software},
    pdfkeywords={blablabla, blablabla, blablabla, blablabla, blablabla},
    colorlinks=false,
    pdfborder={0 0 0}
}
\makeatother

\begin{document}



% ------------------------------------------------------------
% CAPA
% ------------------------------------------------------------
\begin{capa}
    \center

    {\fontsize{12}{14}\selectfont PONTIFÍCIA UNIVERSIDADE CATÓLICA DE CAMPINAS}\par
    {\fontsize{12}{14}\selectfont ENGENHARIA DE SOFTWARE}\par
    {\fontsize{12}{14}\selectfont 12563 - PROJETO INTEGRADOR VI - ENGENHARIA DE SOFTWARE}

    \vfill

    {\fontsize{12}{16}\selectfont
        \textbf{24002529 Eduarda Luiza Pinheiro Nepomuceno}\par
        \textbf{24008214 Jean Yuki Kimura}\par
        \textbf{24014979 Jhenifer Laís Barbosa}\par
        \textbf{24019083 João Pedro Duarte Giatti}\par
    }

    \vfill

    {\bfseries\fontsize{12}{16}\selectfont 
        \imprimirtitulo\par}

    \vfill

    {\fontsize{12}{14}\selectfont
        Orientadora: Profa. Ma. Sílvia Cristina de Matos Soares\par
    }

    \vfill

    {\fontsize{12}{14}\selectfont 
        Campinas\par
        2º Semestre de 2026\par
    }
\end{capa}

% ------------------------------------------------------------
% RESUMO
% ------------------------------------------------------------
\begin{resumo}
    \SingleSpacing
    
    Apresentar uma visão geral do projeto, incluindo o problema abordado, a solução proposta
    com o uso da tecnologia atual (Inteligência Artificial ou Blockchain ou tema p
    Cibersegurança ou Business Intelligence ou outra acordada com a orientadora), a
    metodologia adotada e os principais resultados obtidos.

    \vspace{\onelineskip}

    \noindent
    \textbf{Palavras-chave}: blablabla. blablabla. blablabla. blablabla. blablabla. blablabla.
\end{resumo}

% ------------------------------------------------------------
% SUMÁRIO
% ------------------------------------------------------------
\pdfbookmark[0]{\contentsname}{sumario}
\tableofcontents*
\cleardoublepage
\textual

% ------------------------------------------------------------
% 1 INTRODUÇÃO
% ------------------------------------------------------------
\chapter{Introdução}

Descrever o \textbf{contexto do problema} (qual é o problema que este projeto pretende resolver),
a motivação para o desenvolvimento do projeto, a \textbf{justificativa da solução proposta} (por
que a solução proposta é a melhor, o que ela traz de vantagens e de benefícios?) e o \textbf{escopo
do projeto} (o que o software fará e o que o software não irá fazer).

Elaborar o \textbf{Círculo Dourado.}

% ------------------------------------------------------------
% 2 OBJETIVOS 
% ------------------------------------------------------------
\chapter{Objetivos}

\section{Objetivo Geral}
Descrever o objetivo principal do projeto.

\section{Objetivos Específicos}
Listar os objetivos específicos que contribuem para atingir o objetivo geral.

\begin{alineas}
    \item blablabla blablabla blablabla blablabla blablabla blablabla;
    \item blablabla blablabla blablabla blablabla blablabla blablabla;
    \item blablabla blablabla blablabla blablabla blablabla blablabla.
\end{alineas}

% ------------------------------------------------------------
% 3 Representantes da Comunidade Externa
% ------------------------------------------------------------
\chapter{Representantes da Comunidade Externa}

Caracterizar a extensão (62 horas de extensão) - são as pessoas que participarão da
definição, das entrevistas, da definição dos requisitos, da validação dos protótipos e da
validação final do projeto).

Nome, contato (e-mail) e cargo do profissional/pessoa e o nome e CNPJ da organização
onde ele(a) atua (2 a 5 profissionais). 

% ------------------------------------------------------------
% 4 Fundamentação Teórica
% ------------------------------------------------------------
\chapter{Fundamentação Teórica}

Apresentar os conceitos teóricos relacionados à tecnologia atual e Engenharia de Software
que embasam o projeto.

Incluir a definição das técnicas/ferramentas/algoritmos etc. das tecnologias atuais
utilizadas.

% ------------------------------------------------------------
% 5 Levantamento e Análise de Requisitos
% ------------------------------------------------------------
\chapter{Levantamento e Análise de Requisitos}

\section{Requisitos Funcionais}
Listar e descrever os requisitos funcionais do sistema (Casos de Uso e demais definições –
histórias de usuários e outros).

\section{Requisitos Não Funcionais}
Listar requisitos de desempenho, segurança, usabilidade, escalabilidade, disponibilidade,
entre outros.

\section{Requisitos Específicos para Inteligência Artificial}
Definir métricas de qualidade do modelo (ex.: acurácia mínima, tempo de resposta, entre
outros).

% ------------------------------------------------------------
% 6 Benchmarking 
% ------------------------------------------------------------
\chapter{Benchmarking}

Elaborar uma tabela para comparar as funcionalidades/benefícios do seu software em
relação aos demais, já existentes no mercado (pelo menos 3 software do mercado).

\vspace{1em}

\begin{table}[h!]
\centering
\renewcommand{\arraystretch}{1.6}
\begin{tabular}{|p{3.5cm}|c|c|c|c|}
\hline
Software & \textbf{Seu Software} & \textbf{Software1} & \textbf{Software2} & \textbf{Software3} \\ \hline
Funcionalidade & & & & \\ \hline
\textbf{Func1} & \cmark & \xmark & \cmark & \cmark \\ \hline
\textbf{Func2} & \cmark & \cmark & \cmark & \cmark \\ \hline
\textbf{Func3} & \xmark & \cmark & \xmark & \xmark \\ \hline
\textbf{Func4} & \cmark & \xmark & \cmark & \xmark \\ \hline
\textbf{Func5} & \cmark & \cmark & \cmark & \xmark \\ \hline
\end{tabular}
\end{table}

Descrever qual é o diferencial do seu software em relação aos softwares já existentes.
Apontar o caráter inovador ou atual do projeto.

% ------------------------------------------------------------
% 7 Definição das Personas 
% ------------------------------------------------------------
\chapter{Definição das Personas}

% ------------------------------------------------------------
% 8 Protótipos das Interfaces 
% ------------------------------------------------------------
\chapter{Protótipos das Interfaces}

Print, legível, das interfaces, no arquivo do relatório final .docx e também o link do
software utilizado para o design (ex: link do figma, XD etc.).

Não será aceito apenas o link.

Validação/aceite dos protótipos com as pessoas/profissionais da comunidade externa.

% ------------------------------------------------------------
% 9 {Modelo de Dados 
% ------------------------------------------------------------
\chapter{Modelo de Dados}

MER para banco de dados relacional ou modelo da estrutura de dados para banco de dados
não relacional.


% ------------------------------------------------------------
% 10 Arquitetura e Tecnologias 
% ------------------------------------------------------------
\chapter{Arquitetura e Tecnologias}

Descrever a arquitetura geral do sistema, incluindo o diagrama/modelo arquitetural e a
integração dos componentes de IA e/ou demais tecnologias atuais do projeto.

Descrever as tecnologias, linguagens, frameworks e bibliotecas utilizadas no
desenvolvimento do sistema (frontend, protótipos, backend, banco de dados, cloud etc).

% ------------------------------------------------------------
% 11 Descrição da Solução de Inteligência Artificial ou outra Tecnologia Atual do Projeto
% ------------------------------------------------------------
\chapter{Descrição da Solução Tecnologica Atual do Projeto}

\section{Base de Dados}
Descrever a origem, características, volume e tratamento dos dados utilizados.

\section{Pré-processamento dos Dados (se utilizar IA)}
Explicar as técnicas de limpeza, normalização e preparação dos dados.

\section{Modelos de IA Utilizados (se utilizar IA)}
Descrever os algoritmos, modelos e hiperparâmetros adotados.

\section{Treinamento e Validação (se utilizar IA)}
Explicar o processo de treinamento, validação e teste dos modelos.

\section{Métricas de Avaliação (se utilizar IA)}
Apresentar e justificar as métricas utilizadas para avaliação do desempenho.

\section{Estrutura definida para Blockchain (se utilizar Blockchain)IA)}
Apresentar e justificar a estrutura e seus relacionamentos definidos para Blockchain.

\section{Definir se considera ataque ou defesa e como (se envolver
Cibersegurança}
Apresentar e justificar se o projeto envolve mecanismos para ataque e/ou defesa p/
Cibersegurança.


% ------------------------------------------------------------
% 12 Testes e Validação do Sistema 
% ------------------------------------------------------------
\chapter{Testes e Validação do Sistema}

Apresentar a estratégia de testes de software e os resultados obtidos.

% ------------------------------------------------------------
% 12 Gestão do Projeto 
% ------------------------------------------------------------
\chapter{Gestão do Projeto}

Descrever a metodologia de desenvolvimento adotada (SCRUM), listar o Product Backlog,
cronograma planejado inicial (SPRINTS) com tarefas e responsáveis, cronograma real final
(SPRINTS) com tarefas e responsáveis.

% ------------------------------------------------------------
% 14 Resultados e Discussão 
% ------------------------------------------------------------
\chapter{Resultados e Discussão}

Analisar os resultados obtidos, confrontando-os com os objetivos propostos.

Descrever quais as facilidades e as melhorias que o projeto traz aos usuários, quais os
objetivos foram atendidos e quais objetivos não foram atendidos (e por que não foram
atendidos?).

Fazer as considerações éticas, de privacidade e de impacto social. Discutir possíveis vieses,
impactos sociais, questões éticas e conformidade com legislações de proteção de dados.

% ------------------------------------------------------------
% 15 Conclusão e Trabalhos Futuros 
% ------------------------------------------------------------
\chapter{Conclusão e Trabalhos Futuros}

Apresentar as conclusões do projeto, limitações encontradas e sugestões de melhorias
futuras.

% ------------------------------------------------------------
% 16 Aceite final do projeto 
% ------------------------------------------------------------
\chapter{Aceite final do projeto}

Elaborar um termo de aceite do projeto, pela comunidade externa.

% ------------------------------------------------------------
% REFERÊNCIAS
% ------------------------------------------------------------
\postextual

\chapter{Referências}

% APENAS DE EXEMPLO
\bibitem{springbootdocs2026}
SPRING.
Spring Boot Reference Documentation.
\textbf{Spring Framework Docs}, 2026.
Disponível em: \url{https://docs.spring.io/spring-boot/docs/current/reference/html/}.
Acesso em: 10 ago. 2026.



\end{document}